%% Free Lawvere models and the all-carrier/nonempty-carrier boundary.

\subsection{Representable models and free generators}
\label{sec:p10-lawvere-free-algebras}

Generator sets and algebra carriers below follow the fixed-universe
convention of Section~\ref{sec:p10-lawvere-theories}.

Use the product-oriented Lawvere theories and models of
Definitions~\ref{def:p10-lawvere-theory} and~\ref{def:p10-lawvere-model}.
Let \(\mathbb T\) be such a theory.
Write \(\operatorname{Mod}(\mathbb T)\) for their category, with all
natural transformations as morphisms, and \(U(M)=M(1)\).
For \(n\in\mathbb N\), let
\[
 R_n=\mathbb T(n,-):\mathbb T\longrightarrow\mathsf{Set},\qquad
 R_n(t)(a)=t\circ a.
\]
The comparison \(c_n^M:M(n)\to M(1)^n\) has coordinates
\(M(p_i^n)\), \(0\leq i<n\).
Powers here are sets of finite tuples; in particular the canonical
map \(A\to A^1\) sends \(a\) to the singleton tuple \((a)\).

\begin{hbtheorem}{thm:p10-lawvere-finite-free}{proved=annex:p10.lawvere-free.1}{Finite free models are representable}
Each \(R_n\) is a model of \(\mathbb T\).
The map
\[
 \eta^n:n\longrightarrow U(R_n)=\mathbb T(n,1),
 \qquad \eta^n(i)=p_i^n,
\]
makes it free on the finite set \(n\): for every model \(M\),
restriction to these generators is a bijection
\[
 \operatorname{Nat}(R_n,M)\longrightarrow\mathsf{Set}(n,U(M)),
 \qquad \alpha\longmapsto\alpha_1\circ\eta^n.
\]
It is natural in \(M\), including arbitrary model homomorphisms.
For \(n=0\), \(R_0\) is initial and its carrier is \(\mathbb T(0,1)\).

Precomposition gives a fully faithful functor
\[
 \mathbb T^{\mathrm{op}}\longrightarrow\operatorname{Mod}(\mathbb T),
 \qquad n\longmapsto R_n,
\]
with
\[
 \operatorname{Nat}(R_n,R_m)\cong\mathbb T(m,n).
\]
Precisely, \(\mathbb T^{\mathrm{op}}\) is isomorphic to the category
whose objects are the indexed chosen models \((n,R_n)\) and whose
arrows are all natural transformations between them.
Every model free on a finite set is isomorphic to one of these;
the isomorphism is unique once the correspondence of generators is fixed.
The generator map is not required to be injective.
\end{hbtheorem}

The opposite-theory and representable-free-model descriptions are
\cite[1963 thesis, Chapter~III, \S1, pp.~74--75;
1968 article, \S1, p.~110]{Lawvere2004FunctorialReprint} and
\cite[Remark~11.21 and Corollary~11.22, pp.~102--103]{AdamekRosickyVitale2010Draft}.
The latter locators are for the February 2010 author draft.
The context-category formulation also appears in
\cite[Remark~4.2.3, p.~33]{Pitts1995Draft}.

\subsection{Free models on arbitrary sets}
\label{sec:p10-lawvere-free-adjunction}

For this construction use the all-carrier category
\(\operatorname{Alg}_{\rm all}(\Sigma,E)\) of
Theorem~\ref{thm:p10-lawvere-models}.
An empty carrier is allowed when all operations can be interpreted.
This convention is local to the Lawvere-model treatment;
Definition~\ref{def:p10-term-algebra} and
Theorem~\ref{thm:p10-free-algebra} retain their nonempty-algebra hypotheses.

For a set \(X\), use a tagged copy of its elements as formal variables,
disjoint from \(\Sigma\), and let \(T_\Sigma(X)\) be the term set.
For terms \(s,t\), put
\[
 s\mathrel{\theta_X}t
 \quad\Longleftrightarrow\quad
 \forall A\in\operatorname{Alg}_{\rm all}(\Sigma,E)\
 \forall v:X\to A\quad s^A(v)=t^A(v),
 \qquad Q_X=T_\Sigma(X)/\theta_X.
\]
The quotient operations are induced by term formation.
When there are no terms, both the relation and the quotient are empty.
Each equation of \(E\) is interpreted universally in its declared
finite variable context, as in Definition~\ref{def:p10-lawvere-presentation}.

\begin{hbtheorem}{thm:p10-lawvere-free-adjunction}{proved=annex:p10.lawvere-free.2}{Free models on sets and the underlying-set adjunction}
For every finitary equational presentation \((\Sigma,E)\) and every
set \(X\), \(Q_X\) is an all-carrier \((\Sigma,E)\)-algebra.
The map \(x\mapsto[x]\) makes it free on \(X\).
For \(v:X\to A\), its unique extending homomorphism is
\[
 \bar v:Q_X\longrightarrow A,\qquad \bar v([t])=t^A(v).
\]
The free construction extends to a functor on sets.

For any Lawvere theory \(\mathbb T\), take the recovery presentation
and identity-on-objects theory isomorphism
\[
 \rho:T_{\Sigma_{\mathbb T},E_{\mathbb T}}\longrightarrow\mathbb T
\]
of Theorem~\ref{thm:p10-lawvere-presentation-recovery}.
With \(Q_X\) formed for this presentation, define
\[
 \mathcal F_{\mathbb T}(X)=M_{Q_X}\circ\rho^{-1},
 \qquad \mathcal F_{\mathbb T}(X)(n)=Q_X^n.
\]
Here \(M_{Q_X}\) is the evaluation model of
Theorem~\ref{thm:p10-lawvere-models}.
There are natural bijections
\[
 \operatorname{Nat}(\mathcal F_{\mathbb T}(X),M)
       \cong\mathsf{Set}(X,U(M)),
\]
so \(\mathcal F_{\mathbb T}\dashv U\).
Its generator map is
\[
 \eta_X:X\longrightarrow U\mathcal F_{\mathbb T}(X)=Q_X^1,
 \qquad \eta_X(x)=([x]).
\]
On a finite set \(n\), this free model is isomorphic to \(R_n\)
by the unique isomorphism respecting the displayed generators.
All assertions include \(X=\varnothing\).
\end{hbtheorem}

The term and presentation constructions are
\cite[\S13, pp.~121--125]{AdamekRosickyVitale2010Draft}; free models on
all sets are discussed in Remark~11.21, pp.~102--103.
The proof below verifies the empty term case separately and selects
representatives only for the finitely many variables in one equation.

\subsection{Empty carriers and the nonempty restriction}
\label{sec:p10-lawvere-empty-carriers}

\begin{hbtheorem}{thm:p10-lawvere-empty-carriers}{proved=annex:p10.lawvere-free.3}{The exact empty-carrier boundary}
Every model \(M\) has \(M(0)\) a singleton.
A model with \(M(1)=\varnothing\) exists exactly when
\[
                         \mathbb T(0,1)=\varnothing.
\]
When it exists it is unique up to its unique natural isomorphism.
It can be written as \(M(0)=\{*\}\), \(M(n)=\varnothing\) for \(n>0\),
with the uniquely possible functions on arrows.

For a presentation \((\Sigma,E)\), arrows \(0\to1\) are singleton tuples
of ground-term classes, hence are canonically in bijection with those
classes.  Thus a basic nullary operation excludes empty
carriers, while the absence of all nullary symbols permits the empty
algebra and empty-carrier model.  Equations identify existing terms;
they do not create a ground term when the grammar has none.

The full subcategory of models with \(M(1)\ne\varnothing\) corresponds
to \(\operatorname{Alg}_{\rm ne}(\Sigma,E)\).
Free models on nonempty generator sets lie in this subcategory.
If \(\mathbb T(0,1)\ne\varnothing\), all models have nonempty carriers
and \(R_0\) is also initial there.
If \(\mathbb T(0,1)=\varnothing\), the all-carrier initial model is
excluded; the nonempty subcategory may or may not have an initial object.
\end{hbtheorem}

For the empty signature with no equations, nonempty algebras are
nonempty sets and have no initial object.  Nonempty semigroups also
have no initial object: every such semigroup has two distinct constant
homomorphisms to the two-element left-zero semigroup, whose product is
\(a b=a\).
In contrast, for one unary symbol \(d\) and the equation
\(d(x)=d(y)\), the one-element algebra is initial among nonempty models,
although there are no ground terms.  The empty algebra is initial in
the all-carrier category.
These examples distinguish an expressible nullary term from a value
defined by a positive-arity operation independently of its argument,
as in \cite[1963 thesis, Chapter~II, \S1, p.~64]{Lawvere2004FunctorialReprint}.
They also explain why the all-carrier adjunction cannot simply be
restricted at the zero-generator input without checking its carrier.
\index{Lawvere theory!free models}\index{free algebra!zero generators}
\index{empty carrier}\index{definable constant}
